\chapter*{不定积分每日一题}

\refer{《积分的方法与技巧》金玉明\ 顾新身\ 毛瑞庭\ /\ 编著}

\exercise{求不定积分$I=\int{\frac{\d x}{a^2\sin^2{x}+b^2\cos^2{x}}}$}

\teacher{张晓鹏}

\postdate{2022年12月7日}

\result{\frac{1}{ab}\arctan{\qty(\frac{a}{b}\tan{x})}+C}

\analysis{
    这题不是很难，上下同除$\cos^2{x}$就可以解决了.但如果之前没遇到过这类题目或者没使用过这类技巧的话，会比较难想.


    下面直接来看吧
    \formula{
        \begin{aligned}
            \int{\frac{\d x}{a^2\sin^2{x}+b^2\cos^2{x}}}&=\int{\frac{\frac{\d x}{\cos^2{x}}}{a^2\tan^2{x}+b^2}}\\
            &=\int{\frac{\d\tan{x}}{a^2\tan^2{x}+b^2}}\\
            &\xlongequal{t=\tan{x}}\int{\frac{\d t}{a^2t^2+b^2}}\\
            &=\int{\frac{\d t}{b^2\qty[\qty(\dfrac{a}{b}t)^2+1]}}\\
            &=\frac{1}{ab}\int{\frac{\d\qty(\dfrac{a}{b}t)}{\qty(\dfrac{a}{b}t)^2+1}}\\
            &=\frac{1}{ab}\arctan{\frac{a}{b}t}+C\\
            &=\frac{1}{ab}\arctan{\qty(\frac{a}{b}\tan{x})}+C
        \end{aligned}
    }
}